% =====================================================================
% Response --- Editor
% ---------------------------------------------------------------------
% The editor's decision letter is not in the repository; the quoted passage
% below is taken from the forwarding message.
% Status: general response written 2026-10-22 (proofread pass: typos,
%         parallelism, nested quotes). Point-level editor items would follow
%         as E2, E3, ... if the letter itself carries further comments.
% =====================================================================
\section*{Editor}

\subsection*{E1 --- Main verdict}
\refcom{[As you will see one report is very negative and one is more
positive. In these circumstance we generally reject the paper. However, in
your case, a member of the Board sent me this message: `I believe it is
possible to provide the authors with a second chance, provided that, in
their revised version, they try to seriously engage and address the major
criticisms raised, especially, by reviewer 2. In conclusion, my summary
recommendation is \emph{Major Revision}, being mindful that this may
ultimately lead to a \emph{Rejection} in view of the very critical comments
received.' On the basis of this evaluation I accepted the idea to give you a
chance.]}

\response Many thanks for granting us the opportunity to submit a revised manuscript. We
agree with the reviewers that our original narrative was somewhat ill-focused,
not only because parts of our original framing kept the most important
aspects out of the spotlight, but also because it led us to ignore certain
foundational issues that damaged the credibility of the original text.
Against this backdrop, we decided to rewrite the whole draft, focusing on a
more exhaustive and general treatment of our main object of interest: the
intersection between foundational assumptions of input-output modeling and
more pragmatic applications in the field of transformational studies that rely
on input-output models as a main workhorse. While being more encompassing in
terms of the model space discussed, this fully rewritten new version allows us
to discuss our original question in a more useful and better substantiated
way.

Concretely, the revision restates the comparison as two primitive questions:
the nature of the short run and who bears the cost. It organizes the comparison
as a matrix of five labor-market closures and three financing closures. As an
example of how input-output modeling and transformation research might interact,
and of how closure choices moderate this interaction, we take the decarbonization
of the German residential sector. The underlying shared framework is stated once for the open economy, with imported intermediates and product taxes carried as explicit leaks rather than absorbed into demand. Equivalence results, collected in Appendices A and B, explain where closures coincide and where they part ways. Two extensions probe the limits of the framework: a Kaldor--Verdoorn version in which technology responds to capacity utilization, and a markup version in which profits enter as a second income stream. Five appendices document the equivalences and the scope of the matrix (A), the model, its closures and the formal proofs (B), the calibration and data pipeline (C), the sectoral responses in every closure (D) and the pricing extensions (E).

It remains to be added that we are extremely thankful for the constructive and direct
feedback received from the reviewers, which guided us in rethinking the main
setup of our narrative.

\loc{The rewrite is the manuscript as a whole; the thesis and the framing are stated in Sec.~1; the two primitive questions, the closure matrix and Table~2 in Sec.~3; the results in Secs.~5--6.}

% AG (Oct 2026): location pointers finalized 2026-10-25.
